\documentclass[11pt,a4paper]{article}

\usepackage[margin=2.5cm]{geometry}
\usepackage{amsmath, amssymb, amsthm}
\usepackage{mathtools}
\usepackage{enumitem}
\usepackage[hidelinks]{hyperref}
\usepackage{xcolor}

\newtheorem{theorem}{Theorem}[section]
\newtheorem{lemma}[theorem]{Lemma}
\newtheorem{corollary}[theorem]{Corollary}
\theoremstyle{definition}
\newtheorem{example}{Example}[section]
\newtheorem{exercise}{Exercise}
\theoremstyle{remark}
\newtheorem*{remark}{Remark}

\newcommand{\R}{\mathbb{R}}
\DeclarePairedDelimiter{\abs}{\lvert}{\rvert}

\title{Olympiad Training Notes\\[2pt] \large Session 6: AM--GM and Cauchy--Schwarz}
\author{Regional Mathematics Olympiad Squad}
\date{Autumn 2026}

\begin{document}
\maketitle

\begin{abstract}
These notes cover the two inequalities that solve more olympiad problems than any others: the arithmetic mean--geometric mean inequality and the Cauchy--Schwarz inequality. For each we give a proof, the conditions for equality, and a set of worked examples chosen to show the main tricks: splitting terms, choosing weights, and rewriting a sum so that the Engel form applies. Exercises with hints are at the end. Most of the material is classical; good longer treatments are \cite{steele2004} and \cite{hlp1952}, and \cite{engel1998} has many more problems.
\end{abstract}

\tableofcontents

\section{The AM--GM inequality}

\begin{theorem}[AM--GM]\label{thm:amgm}
For nonnegative reals $a_1, \dots, a_n$,
\begin{equation}\label{eq:amgm}
  \frac{a_1 + a_2 + \cdots + a_n}{n} \;\geq\; \sqrt[n]{a_1 a_2 \cdots a_n},
\end{equation}
with equality if and only if $a_1 = a_2 = \cdots = a_n$.
\end{theorem}

\begin{proof}
We use Cauchy's forward--backward induction.

\emph{The case $n = 2$.} For $a, b \geq 0$,
\begin{align*}
  \frac{a + b}{2} - \sqrt{ab}
    &= \frac{a - 2\sqrt{ab} + b}{2} \\
    &= \frac{\bigl(\sqrt{a} - \sqrt{b}\bigr)^2}{2} \;\geq\; 0,
\end{align*}
with equality exactly when $a = b$.

\emph{Forward step: $n \Rightarrow 2n$.} Suppose \eqref{eq:amgm} holds for $n$ numbers. Given $2n$ numbers, apply the $n$-case to each half and then the $2$-case to the two means:
\begin{align*}
  \frac{1}{2n}\sum_{i=1}^{2n} a_i
    &= \frac{1}{2}\left( \frac{1}{n}\sum_{i=1}^{n} a_i + \frac{1}{n}\sum_{i=n+1}^{2n} a_i \right) \\
    &\geq \frac{1}{2}\left( \sqrt[n]{a_1 \cdots a_n} + \sqrt[n]{a_{n+1} \cdots a_{2n}} \right) \\
    &\geq \sqrt{\sqrt[n]{a_1 \cdots a_n}\,\sqrt[n]{a_{n+1} \cdots a_{2n}}} \\
    &= \sqrt[2n]{a_1 a_2 \cdots a_{2n}}.
\end{align*}
So the inequality holds for every power of $2$.

\emph{Backward step: $n \Rightarrow n - 1$.} Suppose it holds for $n$ numbers, and let $a_1, \dots, a_{n-1} \geq 0$ have mean $m$. Apply the $n$-case to $a_1, \dots, a_{n-1}, m$:
\begin{align*}
  m = \frac{a_1 + \cdots + a_{n-1} + m}{n}
    &\geq \sqrt[n]{a_1 \cdots a_{n-1}\, m}, \\
  \intertext{so raising to the $n$th power and dividing by $m$ (if $m = 0$ there is nothing to prove),}
  m^{n-1} &\geq a_1 \cdots a_{n-1}.
\end{align*}
Every $n$ lies below some power of $2$, so the forward and backward steps together give all $n$. Equality propagates through each step only when all the numbers are equal.
\end{proof}

\begin{theorem}[Weighted AM--GM]\label{thm:wamgm}
Let $a_1, \dots, a_n > 0$ and let $w_1, \dots, w_n > 0$ with $w_1 + \cdots + w_n = 1$. Then
\begin{equation}\label{eq:wamgm}
  w_1 a_1 + w_2 a_2 + \cdots + w_n a_n \;\geq\; a_1^{w_1} a_2^{w_2} \cdots a_n^{w_n}.
\end{equation}
\end{theorem}

\begin{proof}
The function $\ln$ is concave, so by Jensen's inequality
\[
  \ln\Bigl(\sum_{i} w_i a_i\Bigr) \;\geq\; \sum_i w_i \ln a_i = \ln\Bigl(\prod_i a_i^{w_i}\Bigr),
\]
and $\exp$ is increasing.
\end{proof}

\subsection{Worked examples}

\begin{example}\label{ex:minimize}
Find the minimum of $f(x) = x^2 + \dfrac{16}{x}$ for $x > 0$.
\end{example}

\noindent\textit{Solution.} We want the product of the terms to be constant, so we split $16/x$ into two equal pieces before applying AM--GM with $n = 3$:
\begin{align*}
  x^2 + \frac{16}{x}
    &= x^2 + \frac{8}{x} + \frac{8}{x} \\
    &\geq 3\sqrt[3]{x^2 \cdot \frac{8}{x} \cdot \frac{8}{x}} \\
    &= 3\sqrt[3]{64} = 12.
\end{align*}
Equality needs $x^2 = 8/x$, i.e.\ $x = 2$, and indeed $f(2) = 4 + 8 = 12$. The minimum is $12$.

\begin{remark}
Splitting $16/x$ as $\tfrac{16}{x} \cdot 1$ and using the two-term AM--GM gives $x^2 + 16/x \geq 2\sqrt{16x} = 8\sqrt{x}$, which is true but useless because the right side is not constant. Choosing the split so that the variable cancels in the product is the main skill here.
\end{remark}

\begin{example}\label{ex:product8}
Let $a, b, c > 0$. Prove that $(a + b)(b + c)(c + a) \geq 8abc$.
\end{example}

\begin{proof}
Apply the two-term AM--GM to each factor and multiply:
\begin{align*}
  (a + b)(b + c)(c + a)
    &\geq 2\sqrt{ab} \cdot 2\sqrt{bc} \cdot 2\sqrt{ca} \\
    &= 8\sqrt{a^2 b^2 c^2} \\
    &= 8abc.
\end{align*}
Equality holds iff $a = b = c$.
\end{proof}

\begin{example}\label{ex:sum1}
Let $a, b, c > 0$ with $a + b + c = 1$. Prove that
\[
  \left(\frac{1}{a} - 1\right)\left(\frac{1}{b} - 1\right)\left(\frac{1}{c} - 1\right) \geq 8.
\]
\end{example}

\begin{proof}
Using the constraint, $\frac{1}{a} - 1 = \frac{1 - a}{a} = \frac{b + c}{a}$, and similarly for the others. Then
\begin{align*}
  \prod_{\text{cyc}} \left(\frac{1}{a} - 1\right)
    &= \frac{(b + c)(c + a)(a + b)}{abc} \\
    &\geq \frac{8abc}{abc} && \text{by Example~\ref{ex:product8}} \\
    &= 8.
\end{align*}
Equality holds at $a = b = c = \tfrac13$.
\end{proof}

\begin{example}\label{ex:weighted}
Maximize $xy^2z^3$ over positive reals with $x + y + z = 6$.
\end{example}

\noindent\textit{Solution.} The exponents suggest writing $6$ as a sum of six terms: one copy of $x$, two copies of $y/2$, and three copies of $z/3$. By \eqref{eq:amgm} with $n = 6$,
\begin{align*}
  1 = \frac{x + y + z}{6}
    &= \frac{x + \frac{y}{2} + \frac{y}{2} + \frac{z}{3} + \frac{z}{3} + \frac{z}{3}}{6} \\
    &\geq \sqrt[6]{x \left(\frac{y}{2}\right)^{2} \left(\frac{z}{3}\right)^{3}} \\
    &= \sqrt[6]{\frac{xy^2z^3}{108}}.
\end{align*}
Hence $xy^2z^3 \leq 108$, with equality when $x = y/2 = z/3$, that is $(x, y, z) = (1, 2, 3)$. The maximum is $108$.

\section{The Cauchy--Schwarz inequality}

\begin{theorem}[Cauchy--Schwarz]\label{thm:cs}
For real numbers $a_1, \dots, a_n$ and $b_1, \dots, b_n$,
\begin{equation}\label{eq:cs}
  \left(\sum_{i=1}^{n} a_i b_i\right)^{2} \;\leq\; \left(\sum_{i=1}^{n} a_i^2\right)\left(\sum_{i=1}^{n} b_i^2\right),
\end{equation}
with equality if and only if the vectors $(a_i)$ and $(b_i)$ are proportional.
\end{theorem}

\begin{proof}
Write $A = \sum a_i^2$, $B = \sum a_i b_i$, $C = \sum b_i^2$. If $A = 0$ the claim is trivial. Otherwise, for every real $t$,
\begin{align*}
  0 \leq \sum_{i=1}^{n} (a_i t + b_i)^2
    &= t^2 \sum a_i^2 + 2t \sum a_i b_i + \sum b_i^2 \\
    &= A t^2 + 2B t + C.
\end{align*}
A quadratic in $t$ with positive leading coefficient that is never negative has discriminant at most zero:
\[
  (2B)^2 - 4AC \leq 0 \iff B^2 \leq AC.
\]
Equality means the quadratic has a real root $t_0$, so $a_i t_0 + b_i = 0$ for all $i$.
\end{proof}

An algebraic identity gives a second proof and shows exactly how large the gap is.

\begin{lemma}[Lagrange's identity]\label{lem:lagrange}
\begin{equation}\label{eq:lagrange}
  \left(\sum_{i=1}^{n} a_i^2\right)\left(\sum_{i=1}^{n} b_i^2\right) - \left(\sum_{i=1}^{n} a_i b_i\right)^{2}
  = \sum_{1 \leq i < j \leq n} (a_i b_j - a_j b_i)^2.
\end{equation}
\end{lemma}

The most useful form for competitions is the following.

\begin{corollary}[Engel form, also called Titu's lemma]\label{cor:engel}
For reals $x_1, \dots, x_n$ and positive reals $y_1, \dots, y_n$,
\begin{equation}\label{eq:engel}
  \frac{x_1^2}{y_1} + \frac{x_2^2}{y_2} + \cdots + \frac{x_n^2}{y_n} \;\geq\; \frac{(x_1 + x_2 + \cdots + x_n)^2}{y_1 + y_2 + \cdots + y_n}.
\end{equation}
\end{corollary}

\begin{proof}
Apply \eqref{eq:cs} with $a_i = x_i/\sqrt{y_i}$ and $b_i = \sqrt{y_i}$:
\[
  \left(\sum_i x_i\right)^2 = \left(\sum_i \frac{x_i}{\sqrt{y_i}}\cdot\sqrt{y_i}\right)^2
  \leq \left(\sum_i \frac{x_i^2}{y_i}\right)\left(\sum_i y_i\right). \qedhere
\]
\end{proof}

\subsection{Worked examples}

\begin{example}\label{ex:149}
Let $x, y, z > 0$ with $x + y + z = 1$. Find the minimum of $\dfrac{1}{x} + \dfrac{4}{y} + \dfrac{9}{z}$.
\end{example}

\noindent\textit{Solution.} The numerators are squares, so \eqref{eq:engel} applies directly:
\begin{align*}
  \frac{1}{x} + \frac{4}{y} + \frac{9}{z}
    &= \frac{1^2}{x} + \frac{2^2}{y} + \frac{3^2}{z} \\
    &\geq \frac{(1 + 2 + 3)^2}{x + y + z} \\
    &= 36.
\end{align*}
Equality in Corollary~\ref{cor:engel} requires $x_i / y_i$ to be constant, here $1/x = 2/y = 3/z$, so $(x, y, z) = (\tfrac16, \tfrac13, \tfrac12)$, giving $6 + 12 + 18 = 36$. The minimum is $36$.

\begin{example}[Nesbitt's inequality]\label{ex:nesbitt}
For $a, b, c > 0$,
\[
  \frac{a}{b + c} + \frac{b}{c + a} + \frac{c}{a + b} \geq \frac{3}{2}.
\]
\end{example}

\begin{proof}
Multiply each numerator and denominator by the numerator to create squares, then apply \eqref{eq:engel}:
\begin{align*}
  \sum_{\text{cyc}} \frac{a}{b + c}
    &= \sum_{\text{cyc}} \frac{a^2}{a(b + c)} \\
    &\geq \frac{(a + b + c)^2}{a(b + c) + b(c + a) + c(a + b)} \\
    &= \frac{(a + b + c)^2}{2(ab + bc + ca)}.
\end{align*}
It remains to show $(a + b + c)^2 \geq 3(ab + bc + ca)$. Expanding,
\begin{align*}
  (a + b + c)^2 - 3(ab + bc + ca)
    &= a^2 + b^2 + c^2 - ab - bc - ca \\
    &= \tfrac{1}{2}\left[(a - b)^2 + (b - c)^2 + (c - a)^2\right] \;\geq\; 0. \tag{$\ast$}\label{eq:sos}
\end{align*}
Therefore the sum is at least $\frac{3(ab + bc + ca)}{2(ab + bc + ca)} = \frac{3}{2}$, with equality iff $a = b = c$.
\end{proof}

\begin{example}\label{ex:abc1}
Let $a, b, c > 0$ with $abc = 1$. Prove that $a^2 + b^2 + c^2 \geq a + b + c$.
\end{example}

\begin{proof}
Let $s = a + b + c$. Cauchy--Schwarz with $b_i = 1$ gives $3(a^2 + b^2 + c^2) \geq s^2$, and AM--GM gives $s \geq 3\sqrt[3]{abc} = 3$. Combining,
\begin{align*}
  a^2 + b^2 + c^2
    &\geq \frac{s^2}{3} = s \cdot \frac{s}{3} \\
    &\geq s \cdot 1 = a + b + c. \qedhere
\end{align*}
\end{proof}

\section{The chain of means}

Both inequalities fit into a single chain. For positive $a_1, \dots, a_n$,
\begin{equation}\label{eq:means}
  \underbrace{\sqrt{\frac{\sum a_i^2}{n}}}_{\text{QM}}
  \;\geq\; \underbrace{\frac{\sum a_i}{n}}_{\text{AM}}
  \;\geq\; \underbrace{\sqrt[n]{\textstyle\prod a_i}}_{\text{GM}}
  \;\geq\; \underbrace{\frac{n}{\sum 1/a_i}}_{\text{HM}},
\end{equation}
with equality throughout iff all $a_i$ are equal. QM $\geq$ AM is \eqref{eq:cs} with $b_i = 1$; AM $\geq$ GM is Theorem~\ref{thm:amgm}; GM $\geq$ HM is AM--GM applied to $1/a_1, \dots, 1/a_n$. A useful consequence of AM $\geq$ HM is
\[
  (a_1 + \cdots + a_n)\left(\frac{1}{a_1} + \cdots + \frac{1}{a_n}\right) \geq n^2.
\]

\section*{Strategy checklist}

\begin{enumerate}[label=\arabic*.]
  \item Find the equality case first (try $a = b = c$ or a boundary value). Every step of your proof must be tight there, or you will end with an inequality that is true but too weak.
  \item For AM--GM, split terms so that the product is constant (Examples~\ref{ex:minimize} and~\ref{ex:weighted}).
  \item For fractions with a sum in the denominator, try creating squares in the numerator and applying the Engel form (Example~\ref{ex:nesbitt}).
  \item Homogeneous inequalities can be normalized: assume $a + b + c = 1$ or $abc = 1$ if convenient.
  \item When nothing else works, move everything to one side and try to write it as a sum of squares, as in \eqref{eq:sos}.
\end{enumerate}

\section*{Exercises}

\begin{exercise}\label{exr:cyclic}
For $a, b, c > 0$, show $\dfrac{a}{b} + \dfrac{b}{c} + \dfrac{c}{a} \geq 3$.
\end{exercise}

\begin{exercise}\label{exr:quartic}
For real $x, y$, show $x^4 + y^4 + 8 \geq 8xy$.
\end{exercise}

\begin{exercise}\label{exr:harmonic}
For $a, b, c > 0$, show $(a + b + c)\left(\dfrac1a + \dfrac1b + \dfrac1c\right) \geq 9$.
\end{exercise}

\begin{exercise}\label{exr:engel2}
For $a, b, c > 0$, show $\dfrac{a^2}{b + c} + \dfrac{b^2}{c + a} + \dfrac{c^2}{a + b} \geq \dfrac{a + b + c}{2}$.
\end{exercise}

\begin{exercise}\label{exr:vector}
Let $x, y, z$ be reals with $x^2 + y^2 + z^2 = 1$. Find the maximum of $2x - y + 2z$, and where it is attained.
\end{exercise}

\begin{exercise}[Harder]\label{exr:imo}
Let $a, b, c > 0$ with $abc = 1$. Prove that
\[
  \frac{1}{a^3(b + c)} + \frac{1}{b^3(c + a)} + \frac{1}{c^3(a + b)} \geq \frac{3}{2}.
\]
\end{exercise}

\subsection*{Hints}

\begin{enumerate}[label=\textbf{\arabic*.}, leftmargin=*]
  \item Three-term AM--GM; the product of the terms is $1$.
  \item Write $8 = 4 + 4$ and use AM--GM on $x^4, y^4, 4, 4$. Watch the sign of $xy$.
  \item This is AM $\geq$ HM from \eqref{eq:means}.
  \item Apply \eqref{eq:engel} directly.
  \item By \eqref{eq:cs}, $(2x - y + 2z)^2 \leq (4 + 1 + 4)(x^2 + y^2 + z^2) = 9$. The maximum $3$ is reached at $(x, y, z) = (\tfrac23, -\tfrac13, \tfrac23)$.
  \item IMO 1995, Problem 2. Substitute $x = 1/a$, $y = 1/b$, $z = 1/c$, rewrite each term as $\frac{x^2}{y + z}$, and reduce to Exercise~\ref{exr:engel2}.
\end{enumerate}

\bibliographystyle{plain}
\bibliography{07-references}

\end{document}
